\begin{textsample}{NEG Dim 2 – 2012\_c5 – Score -2.00 – 2012\_c5/enade\_2012\_ci\_cont\_22dcf93b02.txt}  \label{ex:f2_neg_071}
CAPÍTULO 5 DISTRIBUIÇÃO DOS CONCEITOS 5.1 PANORAMA NACIONAL DA DISTRIBUIÇÃO DOS CONCEITOS A Tabela 5.1 apresenta a quantidade e distribuição de cursos de Ciências Contábeis participantes do ENADE/2012, por faixa de conceito e Grande Região.
A diferença entre os cursos tabulados neste capítulo e no capítulo 2 são os cursos sem conceito, em princípio, aqueles sem alunos concluintes que participassem da prova ou que tivessem tão somente um aluno concluinte10.
Observando -se os dados da Tabela 5.1, nota -se que, dos 865 cursos participantes, 390 ( 45,1 \% ) classificaram -se com conceito 3, o valor modal.
Este foi também o conceito modal em todas as regiões do Brasil ( entre 42,4 \% no Sudeste a 49,7 \% no Nordeste ).
O conceito 2 foi o segundo mais frequente em nível nacional ( 27,5 \%, correspondendo a 238 cursos ) e o conceito 4, o terceiro ( 20,8 \%, correspondendo a 180 cursos ).
Houve, ainda, 31 cursos ( 3,6 \% ) que receberam conceito 5 e outros 15 cursos que receberam conceito 1. 11 dos 865 cursos de Ciências Contábeis ficaram sem conceito ( SC ).
A região Norte participou com 63 cursos ou 7,3 \% do total nacional.
Nesta região, o conceito 3 foi atribuído a 28 cursos dos 63 participantes, o que equivale a 44,4 \% do total regional.
Apenas um curso ficou sem conceito ( SC ), enquanto nenhum curso foi alocado ao conceito 5.
Os demais foram avaliados com conceito 1 ( quatro cursos, correspondendo a 6,3 \% ), conceito 2 ( 23 cursos, correspondendo a 36,5 \% ) e conceito 4 ( sete cursos, correspondendo a 11,1 \% ). 10 Estes cursos com somente um concluinte foram considerados no capítulo 2.
A região Nordeste participou com 163 cursos ou 18,8 \% do total nacional.
Como já comentado, destes, 81 cursos, 49,7 \% em termos regionais, obtiveram conceito 3, o conceito modal para a região.
Os conceitos 2 e 4 foram atribuídos a, respectivamente, 53 e 20 cursos ( 32,5 \% e 12,3 \% ).
O conceito 5 foi atribuído a dois cursos ( 1,2 \% ).
Nessa região quatro dos cursos ficaram sem conceito ( SC ), enquanto três ( 1,8 \% ) receberam o conceito 1.
Dos 337 cursos participantes da região Sudeste, como já comentado, 143 ( 42,4 \% ) obtiveram conceito 3, o conceito modal.
O conceito 1 foi atribuído a cinco cursos ( 1,5 \% ) e o conceito 2, a 96 ( 28,5 \% ). 76 cursos ( 22,6 \% ) receberam o conceito 4 e 16 outros ( 4,7 \% ) receberam o conceito 5.
Apenas um curso ficou sem conceito ( SC ).
A região Sul também contou com cursos \textbf{distribuídos} em todas as faixas de conceitos.
A predominância do conceito 3 foi de 46,4 \%, correspondentes a 91 dos 196 cursos participantes na região Sul.
O conceito 2 foi atribuído a 27 cursos ( 13,8 \% ) e o conceito 4, a 66 cursos ( 33,7 \% ).
O conceito 1 foi recebido por um único curso ( 0,5 \% ) e o conceito 5, por oito deles ( 4,1 \% ).
Apenas três dos cursos da região Sul ficaram sem conceito ( SC ).
Dos 106 cursos participantes na região Centro-Oeste, 47 ( 44,3 \% ) receberam conceito 3, o conceito modal.
Nesta região, dois cursos ficaram sem conceito ( SC ) e outros dois com conceito 1 ( 1,9 \% cada ).
Os demais cursos foram avaliados com conceito 2 ( 39 cursos, 36,8 \% ), conceito 4 ( 11 cursos, 10,4 \% ) e conceito 5 ( cinco cursos, 4,7 \% ). 5.2 CONCEITOS POR CATEGORIA ADMINISTRATIVA E POR GRANDE REGIÃO A Tabela 5.2 apresenta a distribuição dos cursos participantes do ENADE/2012 de Ciências Contábeis, por Categoria Administrativa, de acordo com os conceitos por eles alcançados, segundo as Grandes Regiões brasileiras.
Dos 865 cursos participantes, 130 ( 15,0 \% ) eram ministrados em instituições públicas e 735 ( 85,0 \% ), em privadas.
De acordo com as informações da Tabela 5.2, em termos nacionais, na obtenção do conceito 5 entre instituições públicas e privadas, das 31 IES que receberam este conceito, nove eram públicas e 22 privadas.
Dos 130 cursos participantes de IES públicas, o conceito 3 foi o valor modal, atribuído a 52 cursos ( 40,0 \% ).
Nesta categoria, quatro dos cursos ficaram sem conceito ( SC ).
Entre os demais cursos de instituições públicas participantes, três obtiveram conceito 1 ( 2,3 \% da categoria ), 20 receberam conceito 2 ( 15,4 \% ) e 42 foram alocados ao conceito 4 ( 32,3 \% ).
Como já comentado, os nove restantes receberam o conceito 5 ( 6,9 \% ).
Na rede privada, o conceito modal também foi 3, com 338 cursos dos 735 da categoria.
Nesta categoria administrativa, sete dos cursos ficaram sem conceito ( SC ).
Entre os demais cursos participantes, 12 receberam conceito 1 e 218, conceito 2.
O conceito 4 foi atribuído a 138 cursos e o conceito 5, como já comentado, a 22 cursos.
Na análise por região, observa -se que, na região Norte, as instituições privadas participaram com 51 cursos ( 81,0 \% do total regional ), dos quais apenas um ficou sem conceito ( SC ).
O conceito modal para as instituições privadas na região foi 3, com 23 cursos, correspondendo a 45,1 \%.
O conceito 2, apresenta frequência próxima do conceito modal, com 21 cursos correspondendo a 41,2 \%.
Dois cursos receberam conceito 1 ( 3,9 \% ).
Quatro cursos foram alocados ao conceito 4 e nenhum recebeu conceito 5.
As instituições públicas participaram com 12 cursos ( 19,0 \% do total regional ), dos quais nenhum ficou sem conceito ( SC ) ou recebeu conceito 5.
O conceito modal para esta Categoria Administrativa, na região Norte, também foi 3, com cinco cursos, correspondendo a 41,7 \%.
Entre os demais cursos, três foram alocados ao conceito 4, dois ao conceito 1 e outros dois ao conceito 2, nesta categoria na região.
Na região Nordeste, a rede privada concentrou 126 dos 163 cursos participantes, equivalentes a 77,3 \% do total da região, a menor proporção entre todas as regiões.
Nesta Categoria Administrativa na região, \textbf{predominaram} os cursos com conceito 3, 61 cursos correspondendo a 48,4 \% dos cursos oferecidos por IES privadas no Nordeste.
Apenas três cursos não receberam conceito ( SC ).
Os demais foram avaliados com conceitos 1 ( dois cursos correspondendo a 1,6 \% ), 2 ( 47 cursos correspondendo a 37,3 \% ), 4 ( 12 cursos correspondendo a 9,5 \% ) e 5 ( um curso correspondendo a 0,8 \% ).
As instituições públicas dessa região participaram com 37 cursos ( 22,7 \% do total da região ), dos quais 20 obtiveram conceito 3, também o valor modal.
Somente um curso ficou sem conceito ( SC ).
Os 17 demais foram avaliados com os conceitos 1 ( um curso ), 2 ( seis cursos ), 4 ( oito cursos ) e 5 ( também um curso ), nesta combinação de região e categoria.
Na região Sudeste, a proporção de cursos da rede privada, 93,2 \%, foi mais elevada do que nas demais regiões brasileiras, correspondendo a 314 dos 337 cursos participantes.
Nesta categoria, na região Sudeste, o conceito modal foi 3 ( 134 cursos ).
Os demais foram avaliados com conceito 1 ( cinco cursos ), conceito 2 ( 91 cursos ), conceito 4 ( 68 cursos ) e conceito 5 ( 16 cursos ).
Nesta combinação de Categoria Administrativa e Grande Região, todos os cursos receberam conceito.
Entre os 23 cursos em instituições públicas na região Sudeste, a categoria modal também foi a que recebeu conceito 3 ( nove cursos ).
Apenas um curso ficou sem conceito ( SC ).
Os demais foram avaliados com os conceitos, 2 ( cinco curso ) e 4 ( oito cursos ).
Na região Sudeste, nesta categoria, nenhum curso ficou com conceito 1 ou 5.
As instituições privadas concentraram 158 dos 196 cursos participantes da região Sul, 80,6 \% do total regional.
Desses, 81 ficaram com conceito 3, o conceito modal.
Apenas três cursos ficaram sem receber conceito ( SC ).
Os demais cursos, nesta combinação de Categoria Administrativa e Grande Região, foram avaliados com conceito 1 ( um curso ), conceito 2 ( 24 cursos ), conceito 4 ( 46 cursos ) e conceito 5 ( três cursos ).
As instituições públicas na região Sul participaram com 38 cursos ( 19,4 \% ), dos quais nenhum ficou sem conceito ( SC ) ou recebeu conceito 1.
O conceito modal foi 4, com 20 cursos.
Além disso, os conceitos 2, 3 e 5 foram recebidos por três, dez e cinco cursos, respectivamente.
Na região Centro-Oeste, 86 dos 106 cursos participantes eram de instituições privadas ( 81,1 \% em termos regionais ).
Destes, 39 concentraram -se no conceito 3, conceito modal.
Nenhum curso ficou sem conceito ( SC ) nesta categoria.
Os demais receberam conceito 1 ( dois cursos ), 2 ( 35 cursos ), conceito 4 ( oito cursos ) e conceito 5 ( outros dois cursos ).
Dos 20 cursos de instituições públicas, oito foram avaliados no conceito 3, conceito modal.
Dois cursos ficaram sem conceito ( SC ).
Os demais receberam conceito 2 ( quatro cursos ), conceito 4 ( três cursos ) e conceito 5 ( outros três cursos ).
Nesta região, nenhum curso de IES pública ficou com conceito 1. 5.3 CONCEITOS POR ORGANIZAÇÃO ACADÊMICA E POR GRANDE REGIÃO Na Tabela 5.3 encontra -se a distribuição dos conceitos atribuídos aos cursos participantes do ENADE/2012 na Área de Ciências Contábeis, por Organização Acadêmica, segundo as Grandes Regiões brasileiras.
Dos 865 cursos participantes, 263 eram oferecidos em Universidades, 103 em Centros Universitários e os demais 499 em Faculdades.
Esta distribuição corresponde a, respectivamente, 30,4 \%, 11,9 \% e 57,7 \% dos cursos.
De acordo com os dados apresentados, de todos os 31 cursos avaliados com conceito 5, 17 eram vinculados a Universidades.
Seis cursos em Universidades ficaram sem conceito ( SC ).
Esse tipo de Organização Acadêmica teve o conceito 3 como modal, com 108 cursos ( 41,1 \% ).
Os demais cursos de Universidades avaliados receberam os conceitos 1 ( cinco cursos ), 2 ( 47 cursos ), 4 ( 80 cursos ) e 5 ( 17 cursos, como já mencionado ).
Entre os cursos em Centros Universitários, o conceito modal também foi 3, com 56 dos 103 cursos ( 54,4 \% ).
Neste tipo de organização acadêmica nenhum curso ficou sem conceito ( SC ).
Os outros cursos neste tipo de Organização Acadêmica receberam os conceitos 1 ( dois cursos ), 2 ( 24 cursos ), 4 ( 18 cursos ) e 5 ( três cursos ).
Nas Faculdades, cinco dos 499 cursos ficaram sem conceito ( SC ) e 226 ( 45,3 \% ) receberam o mesmo conceito modal dos demais tipos de Organização Acadêmica, conceito 3.
Dos demais cursos neste tipo de Organização Acadêmica, oito receberam conceito 1, 167 conceito 2, 82 conceito 4 e 11 cursos, o conceito 5.
Considerando -se separadamente as regiões brasileiras, verifica -se que, na região Norte, as Universidades concentraram 12 dos 63 cursos participantes.
Dos cursos em Universidades, nenhum ficou sem conceito ( SC ) ou recebeu conceito 5.
Somente um curso recebeu conceito 1, dois receberam conceito 2, seis receberam o conceito modal 3 e três receberam o conceito 4.
Os Centros Universitários da região Norte foram representados por seis cursos, nenhum deles sem conceito ( SC ).
Quatro receberam conceito 3, o conceito modal, e dois receberam os conceito 2 e 4 ( um curso em cada ).
As Faculdades participaram com 45 cursos na região Norte.
Apenas um ficou sem conceito ( SC ), enquanto nenhum recebeu o conceito 5.
Os demais receberam conceitos 1 ( três cursos ), 2 ( 20 cursos, conceito modal ), 3 ( 18 cursos ) e 4 ( outros três cursos ).
Na região Nordeste, as Universidades participaram com 45 dos 163 cursos na Área de Ciências Contábeis da região.
Apenas um dos cursos em Universidades no Nordeste ficou sem conceito ( SC ).
O conceito modal foi 3, com 25 cursos.
Os demais receberam conceitos 1 ( um curso ), 2 ( oito cursos ), 4 ( nove cursos ) e 5 ( também um curso ).
Os Centros Universitários contaram com seis cursos participantes na região Nordeste, dois com conceito 2 e quatro com conceito 3, também o conceito modal.
As Faculdades foram representadas por 112 cursos na região Nordeste.
O conceito modal também foi 3, com 52 cursos.
Dos cursos restantes, dois receberam conceito 1, 43 conceito 2, 11 conceito 4 e somente um conceito 5, enquanto três não receberam conceito ( SC ).
Na região Sudeste, as Universidades concentraram 87 dos 337 cursos da região.
Entre os cursos em Universidades na região, o conceito modal foi 3 com 32 cursos, e apenas um curso ficou sem conceito ( SC ).
Os demais cursos receberam os conceitos 1 ( dois cursos ), 2 ( 23 cursos ), 4 ( 24 cursos ) e 5 ( cinco cursos ).
Os Centros Universitários participaram com 66 cursos na região Sudeste, dos quais 34 obtiveram conceito modal, 3, e nenhum ficou sem conceito ( SC ).
Os demais receberam os conceitos 1 ( dois cursos ), 2 ( 17 cursos ), 4 ( 11 cursos ) e conceito 5 ( dois cursos ).
As Faculdades foram representadas por 184 cursos na região Sudeste, que se alocaram nos conceitos 1 ( um curso ), 2 ( 56 cursos ), 3 ( 77 cursos, conceito modal ), 4 ( 41 cursos ) e conceito 5 ( nove cursos ).
Nenhum curso ficou sem conceito ( SC ).
Dos 196 cursos da região Sul, 90 eram de Universidades, para os quais o conceito modal foi 4, com 40 cursos.
Nesse tipo de organização, dois dos cursos ficaram sem conceito ( SC ) e os demais receberam os conceitos 1 ( um curso ), 2 ( cinco cursos ), 3 ( 35 cursos ) e 5 ( sete cursos ).
Os Centros Universitários da região Sul tiveram nove dos 16 cursos participantes no conceito modal, 3.
Os outros sete cursos receberam conceitos 2 ( um curso ), 4 ( cinco cursos ) e 5 ( também um curso ).
Foram 90 os cursos vinculados a Faculdades na região Sul, e apenas um destes ficou sem conceito ( SC ), enquanto nenhum recebeu conceitos 1 ou 5.
Dos demais cursos, 21 receberam conceito 2, 47 conceito 3 ( conceito modal ) e o outros 21, conceito 4.
Na região Centro-Oeste, 29 dos 106 cursos eram de Universidades.
Nesse tipo de organização, dois cursos ficaram sem conceito ( SC ), enquanto nenhum curso recebeu conceito 1.
Além disso, o conceito modal foi 3, com dez cursos.
Os outros cursos obtiveram os conceitos 2 ( nove cursos ), 4 ( quatro cursos ) e 5 ( também quatro cursos ).
Os Centros Universitários da região Centro-Oeste contaram com nove cursos, dos quais três receberam o conceito 2, cinco o conceito 3 e somente um curso, o conceito 4.
Nenhum curso ficou sem conceito ( SC ) ou foi alocado aos conceitos 1 ou 5.
Dos 68 cursos em Faculdades na região Centro-Oeste, nenhum ficou sem conceito ( SC ), dois receberam conceito 1, 27 conceito 2, 32 conceito 3 ( modal ), seis conceito 4 e somente um recebeu conceito 5.

% matched lemmas: distribuir, predominar
\end{textsample}
